\section{The cubic log-gamma moment: a classical evaluation and a source correction}
\label{sec:cubic}

Write
\[
 a=\frac12\log(2\pi),\qquad A=\gamma+\log(2\pi),\qquad
 M_j=\int_0^1\log^j\Gamma(x)\,dx.
\]
The manuscript's statement that identifying a named depth-two atom for $M_3$
is open requires correction.  Bailey, D.~Borwein, and J.~M.~Borwein evaluated
$M_3$ using derivatives of Tornheim--Witten zeta functions in Theorem~6,
equation~(61), of their 2015 paper~\cite{BBB}.  The author-hosted preprint is dated
27 June 2012.  Their formula does not establish a reduction to a prescribed
finite collection of ordinary zeta values and derivatives.  That more
restrictive question should be kept distinct from the already established
Tornheim evaluation.  The argument below independently derives an equivalent
formula and supplies a convenient convergence justification.

For complex parameters in an open neighbourhood of $(1,1,1)$, define
\[
 T(s,t,u)=\sum_{m,n\ge1}\frac1{m^s n^t(m+n)^u},
 \qquad T_j=\partial_jT,\quad T_{ij}=\partial_i\partial_jT.
\]
Every derivative below is evaluated at $(1,1,1)$ unless arguments are shown.
The notation $T_{12}$ and $T_{13}$ is preferable to an unqualified prime,
because different parameter derivatives give different constants.

\begin{theorem}[Classical cubic evaluation, with a direct proof]
Set
\[
 V=\frac{\pi^2}{48}+\frac{A^2}{12}
       -\frac{A\zeta'(2)}{\pi^2}+\frac{\zeta''(2)}{2\pi^2}.
\]
Then
\begin{equation}\label{cubic:main}
 \boxed{\displaystyle
 M_3=a^3+3aV+\frac{3\zeta(3)}{16}
 +\frac{3}{8\pi^2}
   \left(2A^2\zeta(3)-2A T_3+2T_{13}-T_{12}\right).}
\end{equation}
All derivatives on the right are represented by absolutely convergent
double series.  In particular, no analytic continuation or summation
prescription is needed to define these constants.
\end{theorem}

\begin{proof}
Let $f(x)=\log\Gamma(x)$, regarded as an integrable periodic function.
Since $f(x)=-\log x+O(x)$ as $x\downarrow0$ and $f$ is bounded near $1$,
we have $f\in L^p(0,1)$ for every finite $p\ge1$.
Kummer's Fourier expansion gives
\[
 f(x)=a+\sum_{n\ge1}
 \left(\frac{\cos(2\pi nx)}{2n}
       +\frac{A+\log n}{\pi n}\sin(2\pi nx)\right),
 \qquad 0<x<1.
\]
Equivalently, the positive Fourier coefficients are
\[
 h_n=\frac1n\left(\frac14-
                   \frac{i(A+\log n)}{2\pi}\right),
 \qquad h_{-n}=\overline{h_n}.
\]
For $0<r<1$, let $f_r$ be the Poisson mean of $f$, so that the coefficient
$h_n$ is replaced by $r^{|n|}h_n$.  The Fourier series is now absolutely
convergent, and Poisson approximate-identity convergence gives
$f_r\to f$ in $L^3$.  Consequently $\int f_r^j\to\int f^j$ for $j=1,2,3$:
for the cubic case this follows, for example, from
\[
 \|u^3-v^3\|_1\le\|u-v\|_3
       (\|u\|_3^2+\|u\|_3\|v\|_3+\|v\|_3^2).
\]

Write $g=f-a$. Orthogonality first gives
\begin{align*}
 \int_0^1g(x)^2\,dx
 &=2\sum_{n\ge1}|h_n|^2\\
 &=\frac{\zeta(2)}8+
   \frac{A^2\zeta(2)-2A\zeta'(2)+\zeta''(2)}{2\pi^2}=V.
\end{align*}
For the cube, a zero-sum triple of nonzero Fourier frequencies has two
positive frequencies $m,n$ and one negative frequency $-(m+n)$, or the
opposite signs.  Accounting for the three positions of the exceptional
sign and then conjugation gives
\begin{equation}\label{cubic:resonance}
 \int_0^1g_r(x)^3\,dx
 =6\mathop{\rm Re}\sum_{m,n\ge1}
     r^{2(m+n)}h_mh_n\overline{h_{m+n}}.
\end{equation}
The factor is correct also when $m=n$: the ordered double sum accounts
for the interchange of the two positive frequencies only when they differ.

These resonant sums converge absolutely at $r=1$. Indeed,
\[
 \sum_{m=1}^{k-1}\frac1{m(k-m)k}
       =\frac{2H_{k-1}}{k^2},
\]
and $|h_mh_nh_{m+n}|$ is bounded by a constant times
$(1+\log(m+n))^3/[mn(m+n)]$.  The grouped majorant is
$O((1+\log k)^4/k^2)$, which is summable.  Dominated convergence therefore
permits $r\uparrow1$ in \eqref{cubic:resonance}.

Put $p=1/4$ and $q_j=(A+\log j)/(2\pi)$.  Direct multiplication yields
\[
 \mathop{\rm Re}\bigl[(p-iq_m)(p-iq_n)(p+iq_{m+n})\bigr]
 =p^3+p(q_mq_{m+n}+q_nq_{m+n}-q_mq_n).
\]
The numerator in the second term is, apart from $4\pi^2$,
\[
 A^2+2A\log(m+n)+\log m\log(m+n)+\log n\log(m+n)-\log m\log n.
\]
By the definitions of the Tornheim derivatives, symmetry in $m,n$,
and $T(1,1,1)=2\zeta(3)$, we obtain
\[
 \int_0^1g(x)^3\,dx
 =\frac{3\zeta(3)}{16}
 +\frac{3}{8\pi^2}
       \bigl(2A^2\zeta(3)-2AT_3+2T_{13}-T_{12}\bigr).
\]
Finally, $\int g=0$ and $\int f^3=a^3+3a\int g^2+\int g^3$.
\end{proof}

\paragraph{Holomorphy and differentiation.}
For completeness, local holomorphy is also elementary.  If each real part
of $s,t,u$ is at least $1-\epsilon$, with $0<\epsilon<1/3$, grouping by
$k=m+n$ and splitting at $m=k/2$ gives a locally uniform majorant of order
$k^{-2+3\epsilon}$.  The sum converges uniformly on compact subsets of
such a neighbourhood.  Standard termwise differentiation therefore applies.
Alternatively, at $(1,1,1)$ any prescribed logarithmic numerator of degree
$d$ is dominated by $2H_{k-1}(\log k)^d/k^2$ after grouping.  These are
ordinary convergent derivatives; Abel means are used only to justify the
Fourier multiplication.

\paragraph{Relations to Euler double zeta derivatives.}
Define $Z(s)=\zeta(s,1)=\sum_{k\ge2}H_{k-1}k^{-s}$, so that the prime
in $Z'(2)$ differentiates its first argument.  Then
\begin{equation}\label{cubic:firstpartials}
 T(1,1,u)=2Z(u+1),\qquad
 T_3=2Z'(2),\qquad T_1=T_2=Z'(2)+\zeta'(3).
\end{equation}
The first identity follows from $1/[m(k-m)]=(1/m+1/(k-m))/k$.
For the last identity, use
$\sum_{n\ge1}1/[n(m+n)]=H_m/m$.
Equation~\eqref{cubic:main} thus exactly matches the published cubic formula.
It also makes clear why a search involving only ordinary Riemann-zeta
derivatives need not find a relation.  Failure of a finite PSLQ search
does not prove that such a reduction is impossible.

\subsection{A positive summand form and an explicit tail bound}

The combination in \eqref{cubic:main} can be computed without cancellation
between separate Tornheim constants.  For $k=m+n$, define
\[
 R_A(m,n)=A^2+2A\log k+\log m\log k+\log n\log k-\log m\log n.
\]
Since $A>0$, $\log k\ge\log m,\log n\ge0$, we have $R_A(m,n)>0$.
Writing $p=\log m,q=\log n,r=\log k$, one also has
$r(p+q)-pq\le r^2$: for fixed $r$, the expression is increasing separately
in $p,q\in[0,r]$.  Consequently
\[
 0<R_A(m,n)\le(A+\log k)^2.
\]
Therefore the diagonal partial sums
\[
 C_N=\frac3{32}\sum_{m+n\le N}\frac1{mn(m+n)}
    +\frac{3}{8\pi^2}
       \sum_{m+n\le N}\frac{R_A(m,n)}{mn(m+n)}
\]
increase to the centered cubic moment $M_3-a^3-3aV$.  Moreover,
\begin{align*}
 0<(M_3-a^3-3aV)-C_N
 &\le\sum_{k>N}\frac{2(1+\log k)}{k^2}
       \left(\frac3{32}+\frac{3(A+\log k)^2}{8\pi^2}\right).
\end{align*}
This gives a transparent computable enclosure, although its convergence
is only $O(\log^3N/N)$.  An integral bound may be used once the explicit
summand is decreasing; direct summation can handle the finite initial range.
This positivity observation is an immediate corollary of the Fourier
calculation, not a claim of a new classical-constant evaluation.

\subsection{Exponentially convergent numerical reproduction}

For $t>0$, put
\[
 C(t)=\operatorname{Li}_1(e^{-t})=-\log(1-e^{-t}),\qquad
 F(t)=\left.\partial_s\operatorname{Li}_s(e^{-t})\right|_{s=1}.
\]
The usual Mellin representation of the Tornheim sum gives
\begin{equation}\label{cubic:mellinpartials}
 T_{12}=\int_0^\infty F(t)^2\,dt,\quad
 T_{13}=\int_0^\infty F(t)C(t)(\log t+\gamma)\,dt,\quad
 T_3=\int_0^\infty C(t)^2(\log t+\gamma)\,dt.
\end{equation}
For $0<t<2\pi$, cancellation of the two poles in Jonqui\`ere's expansion
produces
\begin{align*}
 C(t)&=-\log t+\sum_{k\ge1}c_kt^k,
 &c_k&=\frac{(-1)^k\zeta(1-k)}{k!},\\
 F(t)&=-\alpha-\gamma\log t-\tfrac12\log^2t+\sum_{k\ge1}b_kt^k,
 &b_k&=\frac{(-1)^k\zeta'(1-k)}{k!},\\
 \alpha&=\gamma_1+\frac{\gamma^2}2+\frac{\pi^2}{12}.
\end{align*}
Here the Stieltjes convention is
$\zeta(1+z)=z^{-1}+\sum_{j\ge0}(-1)^j\gamma_jz^j/j!$.
For computational stability the coefficient formulas can be rewritten as
\begin{align*}
 c_1&=\tfrac12,&b_1&=\tfrac12\log(2\pi),\\
 c_{2j}&=\frac{(-1)^j\zeta(2j)}{j(2\pi)^{2j}},&
 b_{2j}&=c_{2j}\left(\log(2\pi)-\psi(2j)-\frac{\zeta'(2j)}{\zeta(2j)}\right),\\
 c_{2j+1}&=0\quad(j\ge1),&
 b_{2j+1}&=\frac{(-1)^{j+1}\pi\zeta(2j+1)}{(2j+1)(2\pi)^{2j+1}}
 \quad(j\ge1).
\end{align*}
The functional equation of the Riemann zeta function proves these identities.
Elementary bounds $\zeta(k)<2$, $|\zeta'(k)|<2$, $H_{k-1}\le k-1$,
$\gamma<1$, $\log(2\pi)<2$, and $2\pi>6$ give, for $k\ge2$,
\[
 |b_k|\le12\,6^{-k},\qquad |c_k|\le2\,6^{-k}.
\]
Thus after degree $K\ge1$ their remainders on $0<t\le1$ are at most
\[
 \epsilon_b=\frac{12}{5}6^{-K},\qquad
 \epsilon_c=\frac25 6^{-K}.
\]
All integrals of the resulting polynomials in $t$ and $\log t$ are evaluated
using $\int_0^1t^k\log^jt\,dt=(-1)^jj!/(k+1)^{j+1}$.

On $t\ge1$, insert
$C(t)=\sum_{n\ge1}e^{-nt}/n$ and
$F(t)=-\sum_{n\ge1}(\log n)e^{-nt}/n$.  The integral with the logarithmic
factor uses
\[
 \int_1^\infty e^{-kt}(\log t+\gamma)\,dt
 =\frac{E_1(k)+\gamma e^{-k}}k.
\]
Truncate this sum at $m+n\le N$.  A common upper bound for the three
omitted tails is
\[
 B_N=4\sum_{k>N}ke^{-k}
 =\frac{4e^{-(N+1)}((N+1)-Ne^{-1})}{(1-e^{-1})^2}.
\]
For the three integrals in the order $T_{12},T_{13},T_3$, safe total
series-truncation bounds are respectively
\[
 \epsilon_b(12+\epsilon_b)+B_N,\qquad
 5\epsilon_b+15\epsilon_c+2\epsilon_b\epsilon_c+B_N,\qquad
 10\epsilon_c+2\epsilon_c^2+B_N.
\]
For example the elementary estimates
$C(t)\le-\log t+1$ and
$|F(t)|\le4-\log t+\tfrac12\log^2t$ on $(0,1]$ imply all the constants
above.  The sharper bound $\int_0^1|F(t)|dt<6$ follows directly from its
coefficient expansion (one may use $0<\alpha<2$).
An elementary verification of the latter constant bound is available:
the defining limit for $\gamma_1$ and the total variation of
$x\mapsto\log x/x$ on $[1,\infty)$ give $|\gamma_1|\le2/e<3/4$;
combine this with $0<\gamma<3/5$ and $3<\pi<\sqrt{12}$.
This is an analytic truncation certificate; floating-point rounding must
be handled separately to turn an implementation into interval arithmetic.

The independent implementation \texttt{verify\_tornheim.py} evaluates the
three partials and compares \eqref{cubic:main} against direct log-gamma
quadrature.  Its output explicitly labels the rounding qualification.

\subsection{A sign check for an auxiliary identity}

If the June 2012 preprint is used, an auxiliary sign should be checked
before copying its Remark~3.  Put
\[
 S_\gamma=\sum_{m,n\ge1}
        \frac{(\gamma+\log m)(\gamma+\log n)}{mn(m+n)}.
\]
Directly expanding the numerator, with absolutely convergent sums, gives
\[
 T_{12}=S_\gamma+2\gamma\bigl(Z'(2)+\zeta'(3)\bigr)
                  -2\gamma^2\zeta(3).
\]
The coefficient of $S_\gamma$ is positive.  The displayed minus sign in
equations~(63) and~(65) of the author-hosted June 2012 preprint would make
their right side negative, although $T_{12}>0$.  This issue does not
affect the cubic formula proved above.  This note does not assert that
the same typographical error occurs in the final journal typesetting.

